\section*{Erd\H{o}s Problem 469}

\paragraph{Statement and scope.}
A positive integer is semiperfect, or pseudoperfect, if it is a sum of
distinct proper divisors. It is primitive semiperfect if none of its
proper divisors is semiperfect. Does the reciprocal sum over all
primitive semiperfect integers converge?
The current page records an affirmative solution by Zachary J.~Lewis
using GPT 5.6 and Claude Fable 5. This attempt reconstructs useful
exact transition formulas from that preprint, but the global
convergence estimate is unresolved here.

\paragraph{Subset-sum notation.}
For an integer \(n\), let \(\Sigma=\sigma(n)\), \(A=\Sigma-2n\),
and let \(X(n)\), \(Y(n)\) be the subset sums of all divisors and
of proper divisors, respectively, including the empty sum.
Complementing a proper-divisor subset shows
\[
 n\in Y(n)\quad\Longleftrightarrow\quad A\in Y(n),
 \tag{1}
\]
since the total proper-divisor sum is \(n+A\).
Call \(n\) weird if it is abundant and not semiperfect.
For such \(n\), define
\[
 C_n=\{t\geq1:A+t\in Y(n)\},\quad
 g(n)=\min C_n,\quad T(n)=\Sigma/g(n).
\]
The set is nonempty because \(n+A\in Y(n)\); hence \(g(n)\leq n\)
and \(T(n)>2\).

\paragraph{Exact extension by a new prime.}
Let \(p\nmid n\) be prime and \(n\) weird. Coprimality gives
\[
 Y(np)=X(n)+pY(n),\qquad A(np)=pA+\Sigma.
 \tag{2}
\]
Also \(X(n)\) is symmetric under \(x\mapsto\Sigma-x\).
By (1)--(2), \(np\) is semiperfect exactly when
\[
 pt\in X(n)\quad\hbox{for some }t\in C_n.
 \tag{3}
\]
Indeed, \(pA+\Sigma=x+py\) with \(x\in X(n)\), \(y\in Y(n)\)
forces \(y>A\): \(y<A\) would give \(x>\Sigma\), while \(y=A\)
is excluded by weirdness. Write \(y=A+t\); then \(x=\Sigma-pt\),
which belongs to \(X(n)\) exactly when \(pt\) does.

If \(np\) is weird, the same calculation with a positive offset
\(s\) from \(A(np)\) gives the exact gap formula
\[
 g(np)=\min_{\substack{t\in C_n,\ x\in X(n)\\x<pt}}(pt-x).
 \tag{4}
\]
The exclusion of \(y\leq A\) remains valid with positive \(s\).
In particular, if \(p>T(n)\), then \(pg(n)>\Sigma\), so (3)
is impossible. The child is still abundant, by multiplicativity
of \(\sigma(n)/n\), and therefore is weird. Formula (4) then gives
\[
 g(np)=pg(n)-\Sigma,\qquad
 T(np)=\frac{T(n)(p+1)}{p-T(n)}.
 \tag{5}
\]
The minimum is attained at \(t=g(n)\), \(x=\Sigma\).
If instead \(p\leq T(n)\) and the child remains weird, take
\(t=g(n)\), \(x=0\) in (4) to obtain
\[
 g(np)\leq pg(n),\qquad
 T(np)\geq T(n)(1+1/p)>p.
 \tag{6}
\]
Thus a first semiperfect extension at a new prime must use
\(p\leq T(n)\), but missed candidates can increase the next threshold.

\paragraph{A fully checked primitive example beyond deficient parents.}
For \(n=70\), \(\Sigma=144\), \(A=4\), and the proper divisors are
\(1,2,5,7,10,14,35\). They cannot sum to \(4\): using only
\(1,2\) gives at most three, while every other divisor is at least five.
Thus \(70\) is weird, and \(g(70)=1\), since \(5\) is a proper divisor.
Consequently every coprime prime \(p>144\) gives a weird number \(70p\).
For \(p=11\), however, \(11=1+10\in X(70)\), so (3) predicts
semiperfectness, witnessed explicitly by
\[
 770=1+10+11+22+77+110+154+385.
\]
All displayed summands are distinct proper divisors.
The maximal proper divisors of \(770\) are \(385,154,110,70\).
Their divisor sums are \(576,288,216,144\), respectively, so the
first three are deficient and the fourth is weird.
Every divisor of a deficient number is deficient, since
\(\sigma(d)/d\) increases under divisibility.
Every divisor of a non-semiperfect number is non-semiperfect:
semiperfectness passes to multiples by scaling its divisor representation.
It follows that no proper divisor of \(770\) is semiperfect.
Thus \(770\) is primitive semiperfect.

\paragraph{Remaining gap and attribution.}
Equations (3)--(6) give exact local information, not a bound on the
total reciprocal mass of all first semiperfect hits over all prime
extension trees. Finitely many candidate primes at each fixed state
does not imply summability of the entire tree; the varying threshold
and the treatment of roots and repeated prime factors remain essential.
The global mass estimate is not proved here.

Source: Z.~J.~Lewis, \emph{Convergence of the reciprocal sum of primitive
semiperfect numbers}, version 1.2.1, Lemma 2.1,
\url{https://github.com/ZachL111/erdos469-preprint/blob/main/paper/output/pdf/erdos469_v1_2_1_preprint.pdf}.
That version explicitly describes independent human review as pending.
The transition arguments are supplied above and are not taken on
faith from its global claim.
Current problem: \url{https://www.erdosproblems.com/469},
checked 24 September 2026. No local formal-verification claim is made.

\paragraph{Completion estimate.}
Exact prime-extension dynamics and a primitive example are proved; reciprocal-mass control is missing.
\textbf{COMPLETION: 45\%}
